% ============================================================
% CHAPTER 3 — Literature Review (REVISED v3)
% ============================================================

\chapter{Literature Review}
\label{chap:literarture}

This chapter surveys the work relevant to this thesis across four
dimensions: the application of \ac{RL} to process
optimisation~(\S\ref{sec:lit_rl_bpm}), compliance-aware
optimisation~(\S\ref{sec:lit_compliance}), multi-objective optimisation
in \ac{BPM}~(\S\ref{sec:lit_moopt}), and resource-aware
optimisation~(\S\ref{sec:lit_resources}). It closes with a synthesis
table and a precise formulation of the research
gap~(\S\ref{sec:research_gap}).

\section{Process Optimisation using RL}
\label{sec:lit_rl_bpm}
Reinforcement learning has entered business process management mainly as
operational decision support: selecting the next activity to execute,
allocating resources to tasks, or recommending interventions that steer
a running case toward a desired outcome. This section traces that work
from early value-based formulations to more recent deep and
reward-driven methods. Two features recur across it. The underlying
process model is almost always imperative or simulation-based rather
than declarative, and compliance is treated as a statistical tendency
the agent is expected to learn rather than a property the environment
guarantees.

\subsection{RL in Business Process Management}

The earliest systematic application of \ac{RL} to \ac{BPM} is due to
Huang et al.~\cite{Huang2011ReinforcementManagement}, who model the
resource allocation problem as an \ac{MDP} and solve it with
Q-learning. Their framework observes the current process state, given by
pending work items and resource availability, and learns a policy that
minimises total execution cost over time, outperforming heuristics such
as round-robin and shortest-processing-time scheduling.

This work establishes two principles that remain influential: framing
process execution as a sequential decision problem amenable to \ac{RL},
and using the process state as the \ac{MDP} state representation.
However, it operates on an imperative process model, optimises a single
objective, and provides no compliance guarantee beyond what the model
implicitly enforces. These three limitations define the space that the
present thesis addresses.

\subsection{Reward-driven and Deep RL Approaches}

More recent work has adopted deep \ac{RL} for resource allocation.
Meneghello et al.~\cite{Meneghello2024OptimizingLearning} combine hybrid
process simulation with deep \ac{RL}, handling the stochastic dynamics
of real-world processes. Middelhuis et
al.~\cite{Middelhuis2025LearningProcesses} develop rollout-based
algorithms with engineered reward functions, showing that reward design
has a first-order effect on convergence and solution quality. Both,
however, focus on global task--resource assignment over imperative
models, optimise a single objective, and do not enforce compliance;
neither addresses declarative semantics or multi-objective trade-offs.

\section{Compliance-aware Process Optimisation}
\label{sec:lit_compliance}
% ------------------------------------------------------------

Compliance, the requirement that executions respect a formal
specification, has been studied extensively in process mining, but
predominantly as a \emph{verification} task rather than an
\emph{optimisation} objective.

\ac{SMT}-based conformance checking, developed by Felli et
al.~\cite{Felli2023Data-awareSMT, Felli2023Multi-perspectiveApproach},
provides the most rigorous treatment of compliance: alignments between
observed traces and normative models are computed by automated reasoning
over satisfiability modulo theories, with formal correctness guarantees
for data- and time-aware models. Their limitation is scalability: the
underlying problem is NP-hard, and execution times grow rapidly with
process complexity, making them unsuitable for runtime decision-making
or large-scale exploration of the solution space. SAT-based conformance
checking~\cite{Boltenhagen2021OptimizedArtefacts} faces the same
ceiling. Both are best understood as \emph{verification oracles} rather
than optimisation engines.

The observation motivating this thesis follows directly: \ac{SMT} and
\ac{SAT} methods can \emph{verify} whether a trace is compliant, but
they cannot \emph{learn} policies that generate compliant traces
efficiently at scale, while existing \ac{RL} approaches can learn
efficient policies but offer no formal compliance guarantee.

\section{Multi-objective Process Optimisation}
\label{sec:lit_moopt}
Optimising a business process rarely reduces to improving a single
quantity: cost, cycle time, resource effort, and quality typically trade
off against one another. This section reviews how prior work discovers
Pareto-optimal process configurations under such competing objectives,
spanning evolutionary search, prescriptive process monitoring, exact
constraint-based solving, and reward scalarisation. Each approach is
examined for two things: how it represents the cost--duration trade-off,
and whether the configurations it returns are guaranteed to remain valid
executions of the underlying model.

\subsection{Evolutionary and Simulation-based Approaches}

The dominant paradigm for multi-objective optimisation in \ac{BPM} is
the combination of agent-based simulation with evolutionary algorithms,
specifically \ac{NSGA}~II~\cite{Deb2002ANSGA-II}. Kirchdorfer et
al.~\cite{KirchdorferLukasFromPolicies} extend the AgentSimulator
framework with \ac{NSGA}~II to discover Pareto-optimal resource handover
policies, reporting average cost and waiting-time reductions of 37\,\%
and 58\,\% over the as-is process.
Doumeni~\cite{ArtemisDoumeniParetoAgentSimulator} develops this line
further, using \ac{NSGA}~II to generate Pareto-optimal traces from
agent-based simulations and evaluating front quality via hypervolume and
Hausdorff distance.

These approaches share a common architecture: a population of candidate
policies is iteratively evaluated by running full simulations, and
evolutionary operators drive the search. This yields diverse Pareto
fronts but at significant computational cost, since each fitness
evaluation requires one or more complete simulations and population
sizes of hundreds of individuals accumulate substantial wall-clock time.

\subsection{Prediction-based Prescriptive Analytics}
\label{sec:lit_prpa}

A distinct and very recent line approaches multi-objective optimisation
through prescriptive process analytics rather than evolutionary search.
Le et al.~\cite{Le2026BalancingAnalytics} recommend the next
\texttt{(activity, resource)} step of a running case to balance
conflicting KPIs, combining a linearly scalarised utility function over
user-defined weights with a predictive model and a genetic search over
the recommendation space. Their setting is the closest in spirit to the
present work, as both navigate an explicit cost--quality trade-off
through weighted scalarisation, and both observe that intermediate
weights yield balanced improvements between the extremes. Two
differences are central. First, their feasibility module restricts
recommendations to \texttt{(activity, resource)} pairs that were
\emph{empirically observed} in the log, mined as a transition system
over abstracted prefixes; compliance is therefore approximate and
data-dependent, in contrast to the formal, model-derived guarantee
pursued here. The consequence is concrete: to prevent the agent from
recommending a premature \emph{Cancel application}, which shortens time
but voids the outcome, they must exclude that activity from the
prescribable set by hand, whereas a declarative shield rules out the
same action structurally. Second, their output is a per-KPI improvement
curve over the weight, rather than a two-dimensional front of
non-dominated solutions.

\subsection{Pareto-optimal Trace Generation from Declarative Models}

Diaz et al.~\cite{Diaz2024Pareto-OptimalModels} address Pareto-optimal
trace generation directly from declarative models, framing it as a
\ac{COP} over an Extended \ac{DCR} graph. Their Branch-and-Bound
approach produces provably optimal solutions and serves as the exact
baseline for EQ1 in this thesis. It shows that declarative models are a
natural fit for multi-objective trace optimisation, but faces the same
scalability limit as other symbolic methods: it does not scale to large
graphs or long traces.

\subsection{Scalarisation Methods}
\label{sec:lit_scalarisation_methods}

The linear scalarisation strategy used in this thesis, its convex-hull
limitation, and the simplex-lattice construction of the weight grid are
presented in Section~\ref{sec:lit_scalarisation}. The point relevant to
positioning against prior work concerns normalisation. Le et
al.~\cite{Le2026BalancingAnalytics} adopt a more elaborate scheme, mapping
KPIs of different scales to a common range through z-normalisation, a
sigmoid, and min-max scaling; the direct normalisation by per-graph maxima
used in this thesis (Section~\ref{sec:reward_normalised}) is a simpler
alternative suited to the bounded cost and duration values of the benchmarks.

\section{Resource-aware Optimisation}
\label{sec:lit_resources}
% ------------------------------------------------------------

Resource-aware optimisation encompasses approaches that explicitly model
the cost and capability of individual resources within the optimisation
objective. At the scheduling level, classical operations-research
methods treat resource assignment as a variant of job-shop scheduling,
optimising makespan or throughput under capacity constraints; these
scale well but assume a fixed process model and a single objective.

At the \ac{RL} level, resource allocation has most often been addressed
at the \emph{global} or \emph{pool} level: a centralised policy decides,
for each arriving work item, which available resource from a shared pool
to assign. Huang et al.~\cite{Huang2011ReinforcementManagement} encode
resource availability directly in the \ac{MDP} state; Middelhuis et
al.~\cite{Middelhuis2025LearningProcesses} and Meneghello et
al.~\cite{Meneghello2024OptimizingLearning}, discussed in
Section~\ref{sec:lit_rl_bpm}, target the same pool-level assignment
problem. All three operate over imperative or simulation-based process
models, optimise a single objective such as cost or cycle time, and do
not enforce compliance architecturally; none encodes resource profiles
as part of a declarative process specification.

Closer to the granularity adopted in this thesis, Kirchdorfer et
al.~\cite{KirchdorferLukasFromPolicies} model individual resource agents
with differentiated cost and duration profiles and discover
Pareto-optimal handover policies via \ac{NSGA}~II, enabling
\emph{event-level} resource specialisation rather than pool-level
assignment. This is the closest prior work to contribution~C3, but it is
evolutionary rather than learning-based, operates on a simulation model
rather than a declarative one, and provides no compliance guarantee.

The present thesis introduces role-conditioned action spaces within a
declarative \ac{DCR} environment, where each action executes an event
under a specific resource role with differentiated cost and duration
multipliers encoded directly in the process specification. Under
unlimited role availability this choice is separable across events; the
capacity-constrained variant
(Section~\ref{sec:loanapp_capacity}, Chapter~\ref{chap:results}) couples
the per-event role choices into a genuine allocation problem and is shown
to outperform standard allocation heuristics under scarcity, closing the
separability gap left open by an unlimited role budget. No prior work
combines event-level, resource-aware multi-objective optimisation with
architectural compliance enforcement over declarative process models;
this gap is formalised in Section~\ref{sec:research_gap}.

% ------------------------------------------------------------
\section{Research Gap}
\label{sec:research_gap}
% ------------------------------------------------------------

Table~\ref{tab:lit_comparison} positions the most relevant work across
the dimensions that matter for this thesis: declarative semantics,
multi-objective optimisation (MO), formal compliance, resource-awareness
(Res), use of reinforcement learning (RL), a safety or compliance
guarantee \emph{during learning} (Safe), Pareto-front recovery (Pareto),
and scalability to larger problems.


The comparison shows that prior work occupies largely disjoint regions
of this design space. Reinforcement-learning approaches over business
processes~\cite{Huang2011ReinforcementManagement,
Meneghello2024OptimizingLearning, Middelhuis2025LearningProcesses} are
scalable and adaptive, but are defined over imperative models, optimise
a single objective, and give no formal guarantee that generated traces
satisfy process constraints.

\begin{sidewaystable}[p]
\centering
\caption{Comparison of related work across key dimensions.
$\checkmark$ = fully addressed, $\circ$ = partially addressed,
$\times$ = not addressed. Scalability refers to the computational
scalability of the method itself, not to the size of the processes
evaluated.}
\label{tab:lit_comparison}
\normalsize
\renewcommand{\arraystretch}{1.5}
\setlength{\tabcolsep}{6pt}
\begin{tabular}{@{}l c l *{8}{c}@{}}
\toprule
\textbf{Work} & \textbf{Year} & \textbf{Method}
& \textbf{Declarative} & \textbf{MO}
& \textbf{Compliance} & \textbf{Res.}
& \textbf{RL} & \textbf{Safe}
& \textbf{Pareto} & \textbf{Scalability}     \\
\midrule
Huang et al.~\cite{Huang2011ReinforcementManagement}      & 2011 & Q-learning BPM        & $\times$ & $\times$ & $\times$ & $\checkmark$ & $\checkmark$ & $\times$ & $\times$ & $\circ$ \\
Meneghello et al.~\cite{Meneghello2024OptimizingLearning} & 2024 & Deep RL + Simulation  & $\times$ & $\times$ & $\times$ & $\checkmark$ & $\checkmark$ & $\times$ & $\times$ & $\checkmark$ \\
Middelhuis et al.~\cite{Middelhuis2025LearningProcesses}  & 2025 & Rollout-based RL      & $\times$ & $\times$ & $\times$ & $\checkmark$ & $\checkmark$ & $\times$ & $\times$ & $\checkmark$ \\
\addlinespace
Le et al.~\cite{Le2026BalancingAnalytics}                 & 2026 & Multi-KPI PrPA        & $\times$ & $\checkmark$ & $\circ$ & $\checkmark$ & $\times$ & $\times$ & $\circ$ & $\checkmark$ \\
Kirchdorfer et al.~\cite{KirchdorferLukasFromPolicies}    & 2025 & NSGA-II + Simulation  & $\times$ & $\checkmark$ & $\times$ & $\checkmark$ & $\times$ & $\times$ & $\checkmark$ & $\circ$ \\
Doumeni~\cite{ArtemisDoumeniParetoAgentSimulator}         & 2024 & NSGA-II + AgentSim    & $\times$ & $\checkmark$ & $\times$ & $\circ$ & $\times$ & $\times$ & $\checkmark$ & $\circ$ \\
Deb et al.~\cite{Deb2002ANSGA-II}                         & 2002 & NSGA-II               & $\times$ & $\checkmark$ & $\times$ & $\times$ & $\times$ & $\times$ & $\checkmark$ & $\circ$ \\
\addlinespace
Diaz et al.~\cite{Diaz2024Pareto-OptimalModels}           & 2024 & Branch \& Bound (DCR) & $\checkmark$ & $\checkmark$ & $\checkmark$ & $\times$ & $\times$ & $\times$ & $\checkmark$ & $\times$ \\
Felli et al.~\cite{Felli2023Multi-perspectiveApproach}    & 2023 & SMT Conformance       & $\checkmark$ & $\times$ & $\checkmark$ & $\times$ & $\times$ & $\times$ & $\times$ & $\times$ \\
Boltenhagen et al.~\cite{Boltenhagen2021OptimizedArtefacts} & 2021 & SAT Conformance     & $\checkmark$ & $\times$ & $\checkmark$ & $\times$ & $\times$ & $\times$ & $\times$ & $\times$ \\
Groefsema et al.~\cite{Groefsema2024SupportingFamilies}   & 2024 & Declarative BPM Survey& $\checkmark$ & $\times$ & $\checkmark$ & $\times$ & $\times$ & $\times$ & $\times$ & $\circ$ \\
\addlinespace
Roijers et al.~\cite{Roijers2013ADecision-Making}         & 2013 & MOMDP Framework       & $\times$ & $\checkmark$ & $\times$ & $\times$ & $\checkmark$ & $\times$ & $\checkmark$ & $\circ$ \\
Vamplew et al.~\cite{Vamplew2008OnFronts}                 & 2008 & Pareto RL (tabular)   & $\times$ & $\checkmark$ & $\times$ & $\times$ & $\checkmark$ & $\times$ & $\checkmark$ & $\circ$ \\
Abels et al.~\cite{Abels2019DynamicLearning}              & 2019 & Deep MORL             & $\times$ & $\checkmark$ & $\times$ & $\times$ & $\checkmark$ & $\times$ & $\checkmark$ & $\checkmark$ \\
Hayes et al.~\cite{Hayes2022APlanning}                    & 2022 & MORL Survey           & $\times$ & $\checkmark$ & $\times$ & $\times$ & $\checkmark$ & $\times$ & $\checkmark$ & $\checkmark$ \\
\addlinespace
Altman~\cite{AltmanCONSTRAINEDPROCESSES}                  & 1999 & CMDP Theory           & $\times$ & $\circ$ & $\checkmark$ & $\times$ & $\checkmark$ & $\checkmark$ & $\times$ & $\circ$ \\
García \& Fernández~\cite{Garcia2015ALearning}            & 2015 & Safe RL Survey        & $\times$ & $\times$ & $\checkmark$ & $\times$ & $\checkmark$ & $\checkmark$ & $\times$ & $\checkmark$ \\
Alshiekh et al.~\cite{Alshiekh2017SafeShielding}          & 2018 & Shielded RL           & $\times$ & $\times$ & $\checkmark$ & $\times$ & $\checkmark$ & $\checkmark$ & $\times$ & $\checkmark$ \\
Achiam et al.~\cite{Achiam2017ConstrainedOptimization}    & 2017 & CPO                   & $\times$ & $\circ$ & $\checkmark$ & $\times$ & $\checkmark$ & $\checkmark$ & $\times$ & $\checkmark$ \\
\midrule
\textbf{This thesis} & \textbf{--} & \textbf{PPO + DCR Shield + MORL} & $\checkmark$ & $\checkmark$ & $\checkmark$ & $\checkmark$ & $\checkmark$ & $\checkmark$ & $\checkmark$ & $\circ$ \\
\bottomrule
\end{tabular}
\end{sidewaystable}


Symbolic and declarative methods, whether
exact optimisers~\cite{Diaz2024Pareto-OptimalModels} or conformance
checkers~\cite{Felli2023Data-awareSMT, Boltenhagen2021OptimizedArtefacts},
provide strong compliance and optimality guarantees but rely on
combinatorial search that scales poorly. Multi-objective methods,
whether evolutionary~\cite{KirchdorferLukasFromPolicies,
ArtemisDoumeniParetoAgentSimulator} or prediction-based
PrPA~\cite{Le2026BalancingAnalytics}, approximate trade-offs effectively but either require many simulation evaluations or enforce feasibility only
empirically, and none integrates a formal compliance guarantee into a
learning policy. Safe and multi-objective \ac{RL}, finally, supply the
two missing ingredients separately: shielding and constrained
optimisation guarantee safety during
learning~\cite{Alshiekh2017SafeShielding, Garcia2015ALearning,
Achiam2017ConstrainedOptimization}, and \ac{MORL} learns
Pareto-optimal trade-offs~\cite{Roijers2013ADecision-Making,
Abels2019DynamicLearning}, but neither is applied to
declarative process models.

\medskip
\textbf{Key insight.}
No existing framework simultaneously (i) operates on declarative process
models with formal semantics, (ii) guarantees compliance of every trace
generated \emph{during} learning, (iii) supports multi-objective
optimisation over process performance metrics, and (iv) learns policies
rather than relying on exhaustive or evolutionary search. The result is
a structural gap between symbolic correctness and learning-based
scalability as Figure~\ref{fig:lit_taxonomy} shows.


\medskip

\textbf{This thesis addresses the gap by integrating Proximal Policy
Optimisation with Dynamic Condition Response graphs through a shield
intrinsic to the \ac{DCR} execution engine.} Every action proposed by the
agent is checked against the current marking, so only enabled events take
effect and every generated trace is compliant by construction, while
\ac{PPO} learns a policy that optimises scalarised cost--duration
objectives. This is, to our knowledge, the first compliance-by-construction
reinforcement-learning approach to multi-objective optimisation over
declarative process models.

\begin{figure}[p]
\centering
\includegraphics[width=0.7\textwidth]{Pictures/lit_taxonomy.png}
\caption{Taxonomy of existing approaches to reinforcement learning for
business process optimisation. Each cluster summarises a body of work by
its principal strength and limitation; no existing framework occupies the
centre, which combines declarative semantics, shielded learning, and
multi-objective optimisation.}
\label{fig:lit_taxonomy}
\end{figure}
